\chapter{Introdución}
\label{chap:introducion}

\lettrine{G}{alicia} presenta una elevada exposición a emergencias en su territorio, entre las que destacan especialmente los incendios forestales, las inundaciones derivadas de episodios de lluvia intensa, los derrumbes de carreteras o terraplenes y determinados riesgos industriales. En los últimos años han surgido emergencias que ilustran su importancia y su impacto operativo: campañas con incendios que han afectado a superficies significativas según la estadística oficial \cite{EGIFMITECO}, episodios de abril de 2026 que motivaron una decena de incidencias en Ourense y la activación de alertas en las Rías Baixas \cite{AxegaIncidenciasChuviaOurense2026,AxegaAlertaChuvia2026}, y movimientos sísmicos locales perceptibles por la población, como el registrado en Ponteceso en 2026 \cite{AxegaSismoPonteceso2026}. La complejidad operativa de la respuesta exige coordinación entre múltiples actores, toma de decisiones en ventanas temporales reducidas y disponibilidad de información operativa fiable. Estos retos son tanto tecnológicos como organizativos y normativos, ya que confluyen distintos niveles administrativos y tecnologías

Este \acrfull{tfm} se apoya y amplía el trabajo previo desarrollado en el \acrfull{tfg} "Aplicación para coordinación de equipos para la prevención y extinción de incendios forestales" \cite{TFGAdrianGarcia2025}, que documenta el diseño e implementación de una aplicación web destinada a centralizar avisos, recursos y despliegues en el contexto de incendios forestales en Galicia. El \acrshort{tfg} aporta decisiones y soluciones de diseño que se toman como punto de partida en este \acrshort{tfm}.

\section{Determinación de la situación actual}
\label{sec:actual}


La situación actual combina capacidades consolidadas de detección, alerta y despliegue con limitaciones de integración entre sistemas y organizaciones. En el ámbito forestal, existen estadísticas nacionales y autonómicas que permiten caracterizar la magnitud del problema \cite{EGIFMITECO}. Galicia dispone además de una estructura operativa asentada en torno a \acrfull{axega} y cuenta con planes de emergencia por tipología de riesgo a nivel autonómico \cite{AxegaPortal,PlansEmerxenciaGalicia}.

Galicia dispone de una organización y planificación detallada para distintas tipologías de riesgo, por ejemplo, \acrfull{pladiga} para riesgos forestales. La Consellería de Presidencia centraliza catálogos, protocolos y guías municipales, la Consellería de Medio Rural publica documentación técnica y memorias sobre defensa del monte, y \acrshort{axega} mantiene boletines e instrumentos operativos. Sin embargo, estos planes y protocolos están dispersos entre distintos organismos y repositorios, y no están unificados en una solución digital compartida y accesible tanto para los equipos de gestión de emergencias como para la ciudadanía común \cite{PlansEmerxenciaGalicia,PladigaXunta2026,AxegaPortal}.

No obstante, persisten dificultades operativas que limitan la eficacia conjunta. La fragmentación de la información hace que avisos ciudadanos, información meteorológica, mapas de emergencias y estado de recursos residan en sistemas segregados. A ello se suma la heterogeneidad organizativa, ya que organismos y equipos tienen procedimientos y necesidades distintas, desde la coordinación estratégica hasta el mando táctico y la intervención en campo. También destaca la variabilidad del riesgo y la concurrencia de eventos, porque no es realista tratar las emergencias de forma aislada y algunas suelen estar relacionadas con otras, lo que exige un modelo multi-riesgo \cite{PlansEmerxenciaGalicia,AxegaAlertaChuvia2026}. Por último, la presión temporal y la trazabilidad obligan a disponer de actualización en tiempo real y de un registro de actuaciones para análisis posteriores.

El análisis realizado en el \acrshort{tfg} muestra que estas limitaciones no son debidas a la ausencia de medios, sino a falta de cohesión operativa digital. Por ello, este \acrshort{tfm} propone una evolución del sistema desarrollado en el \acrshort{tfg} hacia una plataforma de coordinación que mejore la integración, la trazabilidad y el soporte a la decisión en contextos multi-riesgo.

\section{Alcance y objetivos}
\label{sec:objetivos}

El alcance de este \acrshort{tfm} es evolucionar la solución técnica y los modelos operativos desarrollados en el \acrshort{tfg} hacia una plataforma de coordinación de emergencias multi-riesgo para Galicia, con dos componentes principales: una plataforma web para gestión y coordinación, y una aplicación móvil nativa para uso en campo y participación ciudadana. Este \acrshort{tfm} parte explícitamente del trabajo técnico y de campo documentado en el \acrshort{tfg} \cite{TFGAdrianGarcia2025}.

La solución propuesta se alinea con el trabajo realizado por protección civil y otros organismos de respuesta a emergencias, priorizando la interoperabilidad entre estos organismos, sin pretender sustituir sistemas institucionales existentes, sino complementarlos aportando una capa unificada de gestión, visualización y soporte a la decisión \cite{LeySNPC17_2015,NormaBasicaPC2023}.

\subsection{Objetivos principales}

A continuación se enumeran los principales objetivos que debe cumplir la aplicación a
desarrollar en este trabajo:

\begin{enumerate}
    \item Extender el dominio funcional del \acrshort{tfg} hacia un enfoque multi-riesgo (incendios, inundaciones, derrumbes y riesgos industriales) acorde con los planes autonómicos.
    \item Ampliar funcionalidades de la plataforma web existente que centraliza la gestión de alertas, emergencias, recursos y despliegues operativos.
    \item Diseñar e implementar una aplicación móvil nativa en Kotlin para reporte ciudadano y soporte a la intervención en campo.
    \item Mejorar la coordinación multi-organizativa, integrando en un mismo flujo de trabajo a coordinadores, brigadas, protección civil y ciudadanía.
    \item Reforzar la trazabilidad y el análisis post-emergencia mediante históricos e indicadores operativos.
\end{enumerate}


\subsection{Objetivos concretos}

\begin{enumerate}
    \item Contextualización: esclarecer el valor aportado por el \acrshort{tfm} y definir el nuevo dominio multi-riesgo.
    \item Definir la arquitectura de la aplicación móvil y entregar un esqueleto que permita añadir funcionalidades progresivamente.
    \item Introducir un modelo unificado de recursos y organizaciones y definir sus capacidades y área operativa.
    \item Permitir la creación y edición de emergencias (por punto o por cuadrante) y la creación o propuesta manual de asignaciones.
    \item Introducir un registro de logs operativos de las emergencias y de las asignaciones de recursos que permita el análisis posterior de las emergencias y la generación de indicadores. 
    \item Integrar notificaciones móviles mediante \acrfull{fcm}, registrar dispositivos y definir el flujo de notificación y aceptación de asignaciones.
    \item Implementar el almacenamiento de reglas basándose en los protocolos actuales y un motor de reglas que sea capaz de generar propuestas de asignación automáticamente.
\end{enumerate}


\section{Estructura de la memoria}
\label{sec:estructura}

La memoria se organiza en once capítulos y anexos. A continuación se ofrece un resumen del contenido de cada capítulo:

\begin{enumerate}
    \item Capítulo 1: Presenta el contexto general, describe la situación actual en Galicia, expone las motivaciones del trabajo y enuncia los objetivos generales y concretos del proyecto.
    \item Capítulo \ref{chap:basetecnoloxica}: Detalla la base tecnológica empleada: componentes de backend y frontend, bases de datos, servicios y librerías relevantes, y las decisiones arquitectónicas que sustentan la implementación.
    \item Capítulo \ref{chap:estadoarte}: Revisa el estado del arte, compara soluciones existentes y extrae lecciones aplicables al diseño de la plataforma propuesta.
    \item Capítulo \ref{chap:viabilidade}: Analiza la viabilidad del proyecto desde las perspectivas técnica, organizativa y normativa, e incluye una estimación de recursos y riesgos.
    \item Capítulo \ref{chap:desenrolo}: Documenta el desarrollo realizado, describiendo la arquitectura global del sistema, los módulos principales y las integraciones implementadas.
    \item Capítulo \ref{chap:planificacionycostes}: Presenta la planificación temporal del trabajo, los hitos y la metodología seguida para organizar las iteraciones de desarrollo.
    \item Capítulo \ref{chap:requisitos}: Especifica los requisitos funcionales y no funcionales que guían el diseño y la implementación de la plataforma.
    \item Capítulo \ref{chap:deseño}: Describe el diseño funcional y técnico, incluyendo el modelo de datos, las APIs, los flujos de interacción y las decisiones de UX/UI.
    \item Capítulo \ref{chap:implementacion}: Detalla la implementación: estructura de código, despliegue, configuraciones y ejemplos de componentes desarrollados.
    \item Capítulo \ref{chap:probas}: Documenta las pruebas realizadas (unitarias, de integración y de aceptación), su cobertura y los resultados obtenidos.
    \item Capítulo \ref{chap:conclusions}: Resume las conclusiones, las aportaciones del trabajo y propone líneas futuras de mejora y continuidad.
\end{enumerate}
